% =============================================================================
%  Section: Experimental Setup
%  H1 Charged-Hadron SIDIS High-Q^2 Analysis Note
% =============================================================================
\section{Experimental Setup}
\label{sec:experimental_setup}

The measurement uses $\ep$ collisions delivered by HERA and recorded with the
H1 detector.  The center-of-mass energy is $\rts=319~\GeV$.  Charged tracks are
reconstructed in the H1 central tracking system. The scattered lepton and the
hadronic final state are reconstructed with the standard H1 high-$\Qsq$ DIS
selection and calibration chain. For more details on the exact algorithms used for high-$Q^2$ analyses, see Refs.~\cite{Kogler:2011zz,Britzger:2013,Peez:2005Hadroo2}

The charged-SIDIS data sample corresponds to an integrated luminosity of
$351~\pbinv$, collected from 2003-2007. This is the luminosity used for the high-$\Qsq$ charged
hadron measurement documented here and for the corresponding low-$\Qsq$ charged
SIDIS analysis.  The luminosity uncertainty included in the high-$\Qsq$ systematic uncertainty budget is a flat $2.7\%$ normalization component.  It is fully
correlated across bins and charges for cross-section-like quantities and cancels
in charge ratios when the same data sample appears in numerator and denominator.

The analysis uses the standard H1 neutral-current DIS event-quality selection,
including the LAr electron acceptance, the S67 high-$\Qsq$ trigger path for data,
and the standard non-ep-background rejection. The exact event selection logic and the
analysis-level fiducial requirements are summarized in Sec.~\ref{sec:selection}.
